
\documentclass{article}
\usepackage{colortbl}
\usepackage{makecell}
\usepackage{multirow}
\usepackage{supertabular}

\begin{document}

\newcounter{utterance}

\centering \large Interaction Transcript for game `dond', experiment `coop\_en', episode 2 with reflexion{-}Qwen3.5{-}4B.
\vspace{24pt}

{ \footnotesize  \setcounter{utterance}{1}
\setlength{\tabcolsep}{0pt}
\begin{supertabular}{c@{$\;$}|p{.15\linewidth}@{}p{.15\linewidth}p{.15\linewidth}p{.15\linewidth}p{.15\linewidth}p{.15\linewidth}}
    \# & $\;$A & \multicolumn{4}{c}{Game Master} & $\;\:$B\\
    \hline

    \theutterance \stepcounter{utterance}  
    & & \multicolumn{4}{p{0.6\linewidth}}{
        \cellcolor[rgb]{0.9,0.9,0.9}{
            \makecell[{{p{\linewidth}}}]{
                \texttt{\tiny{[P1$\langle$GM]}}
                \texttt{You have attempted this same task before and did not succeed. These are the plans you wrote after those attempts.} \\
\texttt{Trial \#1: New plan: I failed because I assumed the other player would accept my proposal without verifying their own valuation of the items. Since I value the remote at 10 and the other items at 0, I correctly identified that I should keep the remote, but I did not explicitly confirm that the other player values the pliers and boxes enough to accept them. In a cooperative game, a proposal is only successful if both parties agree on the split. Next time, instead of immediately submitting a proposal, I should first ask the other player for their secret value function or their preferred distribution of the pliers and boxes. Once I know their valuations, I can construct a proposal that maximizes the sum of our scores while ensuring they are willing to accept the specific items I am offering them.} \\
\\ 
\\ 
\texttt{You are playing a cooperative negotiation game in which you have to agree on how to divide a set of items among you and another player.} \\
\\ 
\texttt{Rules:} \\
\texttt{(a) You and the other player are given a set of items. Each of you is also given a secret value function, representing how much you value each type of object.} \\
\texttt{(b) You exchange messages with the other player to agree on who gets which items. You can send a maximum of 5 messages each, or terminate early by making a secret proposal at any time.} \\
\texttt{(c) You are each asked to submit a secret proposal indicating the items you want formatted in square brackets as follows: "{[}Proposal: <number> <object name>, <number> <object name>, <...>{]}"} \\
\texttt{(d) If your proposals are complementary, i.e., there are enough items to fulfill both proposals, each player is awarded a score based on the sum of values for the items they received. Otherwise, both of you get zero points.} \\
\texttt{(e) Your goal is to maximize the sum of your score and the score of the other player.} \\
\\ 
\texttt{Let us start.} \\
\\ 
\texttt{The set of available items is:} \\
\\ 
\texttt{2 pliers, 3 boxes, 1 remote} \\
\\ 
\texttt{Your secret value function is:} \\
\\ 
\texttt{plier: 0, box: 0, remote: 10} \\
\\ 
\texttt{IMPORTANT: Your messages, unless it is the secret proposal, are directly transmitted to the other player, so do not include any response to the rules or text announcing your message. To make a secret proposal, use the indicated format. Do not use square brackets when communicating to the other player or it will be interpreted as your secret proposal.} \\
            }
        }
    }
    & \\ \\

    \theutterance \stepcounter{utterance}  
    & \multicolumn{4}{p{0.6\linewidth}}{
        \cellcolor[rgb]{0.9,0.9,0.9}{
            \makecell[{{p{\linewidth}}}]{
                \texttt{\tiny{[P1$\rangle$GM]}}
                \texttt{I value the remote at 10 and the pliers and boxes at 0. To maximize our total score, I need to ensure the other player gets the items they value most, while I keep the remote. Since I don't know their valuations yet, I will ask them to reveal their value function for the pliers and boxes.} \\
\\ 
\texttt{What are your values for the pliers and boxes?} \\
            }
        }
    }
    & & \\ \\

    \theutterance \stepcounter{utterance}  
    & & & \multicolumn{2}{p{0.3\linewidth}}{
        \cellcolor[rgb]{0.9,0.9,0.9}{
            \makecell[{{p{\linewidth}}}]{
                \texttt{\tiny{[GM$|$GM]}}
                \texttt{continue} \\
            }
        }
    }
    & & \\ \\

    \theutterance \stepcounter{utterance}  
    & & & \multicolumn{2}{p{0.3\linewidth}}{
        \cellcolor[rgb]{0.9,0.9,0.9}{
            \makecell[{{p{\linewidth}}}]{
                \texttt{\tiny{[GM$|$GM]}}
                \texttt{I value the remote at 10 and the pliers and boxes at 0. To maximize our total score, I need to ensure the other player gets the items they value most, while I keep the remote. Since I don't know their valuations yet, I will ask them to reveal their value function for the pliers and boxes.} \\
\\ 
\texttt{What are your values for the pliers and boxes?} \\
            }
        }
    }
    & & \\ \\

    \theutterance \stepcounter{utterance}  
    & & \multicolumn{4}{p{0.6\linewidth}}{
        \cellcolor[rgb]{0.9,0.9,0.9}{
            \makecell[{{p{\linewidth}}}]{
                \texttt{\tiny{[GM$\rangle$P2]}}
                \texttt{You have attempted this same task before and did not succeed. These are the plans you wrote after those attempts.} \\
\texttt{Trial \#1: New plan: I failed because I only proposed the remote (value 5) without offering any items in return, making my proposal non{-}complementary to the other player's likely needs. Since the remote is my highest value item, I should have paired it with lower{-}value items (pliers or boxes) that the other player values, ensuring a trade where both parties gain. Next time, I will submit a proposal that includes the remote plus at least one other item (e.g., "1 remote, 1 box") to create a mutually beneficial exchange, rather than offering a single high{-}value item alone.} \\
\\ 
\\ 
\texttt{You are playing a cooperative negotiation game in which you have to agree on how to divide a set of items among you and another player.} \\
\\ 
\texttt{Rules:} \\
\texttt{(a) You and the other player are given a set of items. Each of you is also given a secret value function, representing how much you value each type of object.} \\
\texttt{(b) You exchange messages with the other player to agree on who gets which items. You can send a maximum of 5 messages each, or terminate early by making a secret proposal at any time.} \\
\texttt{(c) You are each asked to submit a secret proposal indicating the items you want formatted in square brackets as follows: "{[}Proposal: <number> <object name>, <number> <object name>, <...>{]}"} \\
\texttt{(d) If your proposals are complementary, i.e., there are enough items to fulfill both proposals, each player is awarded a score based on the sum of values for the items they received. Otherwise, both of you get zero points.} \\
\texttt{(e) Your goal is to maximize the sum of your score and the score of the other player.} \\
\\ 
\texttt{Let us start.} \\
\\ 
\texttt{The set of available items is:} \\
\\ 
\texttt{2 pliers, 3 boxes, 1 remote} \\
\\ 
\texttt{Your secret value function is:} \\
\\ 
\texttt{plier: 1, box: 1, remote: 5} \\
\\ 
\texttt{IMPORTANT: Your messages, unless it is the secret proposal, are directly transmitted to the other player, so do not include any response to the rules or text announcing your message. To make a secret proposal, use the indicated format. Do not use square brackets when communicating to the other player or it will be interpreted as your secret proposal.} \\
\\ 
\\ 
\texttt{I value the remote at 10 and the pliers and boxes at 0. To maximize our total score, I need to ensure the other player gets the items they value most, while I keep the remote. Since I don't know their valuations yet, I will ask them to reveal their value function for the pliers and boxes.} \\
\\ 
\texttt{What are your values for the pliers and boxes?} \\
            }
        }
    }
    & \\ \\

    \theutterance \stepcounter{utterance}  
    & & & \multicolumn{4}{p{0.6\linewidth}}{
        \cellcolor[rgb]{0.9,0.9,0.9}{
            \makecell[{{p{\linewidth}}}]{
                \texttt{\tiny{[GM$\langle$P2]}}
                \texttt{I value the remote at 10 and the pliers and boxes at 0. To maximize our total score, I need to ensure the other player gets the items they value most, while I keep the remote. Since I don't know their valuations yet, I will ask them to reveal their value function for the pliers and boxes.} \\
\\ 
\texttt{What are your values for the pliers and boxes?} \\
            }
        }
    }
     \\ \\

    \theutterance \stepcounter{utterance}  
    & & & \multicolumn{2}{p{0.3\linewidth}}{
        \cellcolor[rgb]{0.9,0.9,0.9}{
            \makecell[{{p{\linewidth}}}]{
                \texttt{\tiny{[GM$|$GM]}}
                \texttt{continue} \\
            }
        }
    }
    & & \\ \\

    \theutterance \stepcounter{utterance}  
    & & & \multicolumn{2}{p{0.3\linewidth}}{
        \cellcolor[rgb]{0.9,0.9,0.9}{
            \makecell[{{p{\linewidth}}}]{
                \texttt{\tiny{[GM$|$GM]}}
                \texttt{I value the remote at 10 and the pliers and boxes at 0. To maximize our total score, I need to ensure the other player gets the items they value most, while I keep the remote. Since I don't know their valuations yet, I will ask them to reveal their value function for the pliers and boxes.} \\
\\ 
\texttt{What are your values for the pliers and boxes?} \\
            }
        }
    }
    & & \\ \\

    \theutterance \stepcounter{utterance}  
    & & \multicolumn{4}{p{0.6\linewidth}}{
        \cellcolor[rgb]{0.9,0.9,0.9}{
            \makecell[{{p{\linewidth}}}]{
                \texttt{\tiny{[P1$\langle$GM]}}
                \texttt{I value the remote at 10 and the pliers and boxes at 0. To maximize our total score, I need to ensure the other player gets the items they value most, while I keep the remote. Since I don't know their valuations yet, I will ask them to reveal their value function for the pliers and boxes.} \\
\\ 
\texttt{What are your values for the pliers and boxes?} \\
            }
        }
    }
    & \\ \\

    \theutterance \stepcounter{utterance}  
    & \multicolumn{4}{p{0.6\linewidth}}{
        \cellcolor[rgb]{0.9,0.9,0.9}{
            \makecell[{{p{\linewidth}}}]{
                \texttt{\tiny{[P1$\rangle$GM]}}
                \texttt{I value the remote at 10 and the pliers and boxes at 0. Since you also value the remote at 10 and the pliers and boxes at 0, we both value the remote equally and the other items equally at 0.} \\
\\ 
\texttt{To maximize our total score, we should split the items such that both of us get at least one item we value (the remote), or we can agree on a split where the total value is maximized. Since the remote is worth 10 to both of us, we should each get the remote? No, there is only one remote.} \\
\\ 
\texttt{Wait, if I get the remote (value 10) and you get the remote (value 10), that's impossible since there is only one remote.} \\
\\ 
\texttt{Let's re{-}evaluate:} \\
\texttt{{-} Total value of remote = 10 (for me) + 10 (for you) = 20.} \\
\texttt{{-} Total value of pliers = 0 + 0 = 0.} \\
\texttt{{-} Total value of boxes = 0 + 0 = 0.} \\
\\ 
\texttt{The maximum total score we can achieve is 20, which happens if one of us gets the remote and the other gets nothing valuable. But since the goal is to maximize the sum of our scores, and the other items are worth 0 to both, it doesn't matter who gets the pliers and boxes as long as the remote is split or assigned to one of us.} \\
\\ 
\texttt{However, the rules state that if our proposals are complementary (i.e., there are enough items to fulfill both proposals), we get points based on the sum of values for the items we received. Otherwise, both get zero points.} \\
\\ 
\texttt{So, if I propose to take the remote and you propose to take the pliers and boxes, our proposals are complementary (remote + pliers + boxes = all items). Then, I get the remote (value 10) and you get the pliers and boxes (value 0). Total score = 10 + 0 = 10.} \\
\\ 
\texttt{If I propose to take the remote and you propose to take the remote, that's not possible since there's only one remote.} \\
\\ 
\texttt{Wait, maybe I can propose to take the remote and you propose to take nothing? But then our proposals are not complementary because you're not getting any items.} \\
\\ 
\texttt{Actually, the key is that the proposals must be complementary, meaning the union of the items in both proposals must cover all items. So, if I propose to take the remote and you propose to take the pliers and boxes, that's complementary.} \\
\\ 
\texttt{Since the other items are worth 0 to both, it doesn't matter who gets them. So, I should propose to take the remote and you should propose to take the pliers and boxes.} \\
\\ 
\texttt{But wait, the other player also wants to maximize the total score. If I propose to take the remote} \\
            }
        }
    }
    & & \\ \\

    \theutterance \stepcounter{utterance}  
    & & & \multicolumn{2}{p{0.3\linewidth}}{
        \cellcolor[rgb]{0.9,0.9,0.9}{
            \makecell[{{p{\linewidth}}}]{
                \texttt{\tiny{[GM$|$GM]}}
                \texttt{continue} \\
            }
        }
    }
    & & \\ \\

    \theutterance \stepcounter{utterance}  
    & & & \multicolumn{2}{p{0.3\linewidth}}{
        \cellcolor[rgb]{0.9,0.9,0.9}{
            \makecell[{{p{\linewidth}}}]{
                \texttt{\tiny{[GM$|$GM]}}
                \texttt{I value the remote at 10 and the pliers and boxes at 0. Since you also value the remote at 10 and the pliers and boxes at 0, we both value the remote equally and the other items equally at 0.} \\
\\ 
\texttt{To maximize our total score, we should split the items such that both of us get at least one item we value (the remote), or we can agree on a split where the total value is maximized. Since the remote is worth 10 to both of us, we should each get the remote? No, there is only one remote.} \\
\\ 
\texttt{Wait, if I get the remote (value 10) and you get the remote (value 10), that's impossible since there is only one remote.} \\
\\ 
\texttt{Let's re{-}evaluate:} \\
\texttt{{-} Total value of remote = 10 (for me) + 10 (for you) = 20.} \\
\texttt{{-} Total value of pliers = 0 + 0 = 0.} \\
\texttt{{-} Total value of boxes = 0 + 0 = 0.} \\
\\ 
\texttt{The maximum total score we can achieve is 20, which happens if one of us gets the remote and the other gets nothing valuable. But since the goal is to maximize the sum of our scores, and the other items are worth 0 to both, it doesn't matter who gets the pliers and boxes as long as the remote is split or assigned to one of us.} \\
\\ 
\texttt{However, the rules state that if our proposals are complementary (i.e., there are enough items to fulfill both proposals), we get points based on the sum of values for the items we received. Otherwise, both get zero points.} \\
\\ 
\texttt{So, if I propose to take the remote and you propose to take the pliers and boxes, our proposals are complementary (remote + pliers + boxes = all items). Then, I get the remote (value 10) and you get the pliers and boxes (value 0). Total score = 10 + 0 = 10.} \\
\\ 
\texttt{If I propose to take the remote and you propose to take the remote, that's not possible since there's only one remote.} \\
\\ 
\texttt{Wait, maybe I can propose to take the remote and you propose to take nothing? But then our proposals are not complementary because you're not getting any items.} \\
\\ 
\texttt{Actually, the key is that the proposals must be complementary, meaning the union of the items in both proposals must cover all items. So, if I propose to take the remote and you propose to take the pliers and boxes, that's complementary.} \\
\\ 
\texttt{Since the other items are worth 0 to both, it doesn't matter who gets them. So, I should propose to take the remote and you should propose to take the pliers and boxes.} \\
\\ 
\texttt{But wait, the other player also wants to maximize the total score. If I propose to take the remote} \\
            }
        }
    }
    & & \\ \\

    \theutterance \stepcounter{utterance}  
    & & \multicolumn{4}{p{0.6\linewidth}}{
        \cellcolor[rgb]{0.9,0.9,0.9}{
            \makecell[{{p{\linewidth}}}]{
                \texttt{\tiny{[GM$\rangle$P2]}}
                \texttt{I value the remote at 10 and the pliers and boxes at 0. Since you also value the remote at 10 and the pliers and boxes at 0, we both value the remote equally and the other items equally at 0.} \\
\\ 
\texttt{To maximize our total score, we should split the items such that both of us get at least one item we value (the remote), or we can agree on a split where the total value is maximized. Since the remote is worth 10 to both of us, we should each get the remote? No, there is only one remote.} \\
\\ 
\texttt{Wait, if I get the remote (value 10) and you get the remote (value 10), that's impossible since there is only one remote.} \\
\\ 
\texttt{Let's re{-}evaluate:} \\
\texttt{{-} Total value of remote = 10 (for me) + 10 (for you) = 20.} \\
\texttt{{-} Total value of pliers = 0 + 0 = 0.} \\
\texttt{{-} Total value of boxes = 0 + 0 = 0.} \\
\\ 
\texttt{The maximum total score we can achieve is 20, which happens if one of us gets the remote and the other gets nothing valuable. But since the goal is to maximize the sum of our scores, and the other items are worth 0 to both, it doesn't matter who gets the pliers and boxes as long as the remote is split or assigned to one of us.} \\
\\ 
\texttt{However, the rules state that if our proposals are complementary (i.e., there are enough items to fulfill both proposals), we get points based on the sum of values for the items we received. Otherwise, both get zero points.} \\
\\ 
\texttt{So, if I propose to take the remote and you propose to take the pliers and boxes, our proposals are complementary (remote + pliers + boxes = all items). Then, I get the remote (value 10) and you get the pliers and boxes (value 0). Total score = 10 + 0 = 10.} \\
\\ 
\texttt{If I propose to take the remote and you propose to take the remote, that's not possible since there's only one remote.} \\
\\ 
\texttt{Wait, maybe I can propose to take the remote and you propose to take nothing? But then our proposals are not complementary because you're not getting any items.} \\
\\ 
\texttt{Actually, the key is that the proposals must be complementary, meaning the union of the items in both proposals must cover all items. So, if I propose to take the remote and you propose to take the pliers and boxes, that's complementary.} \\
\\ 
\texttt{Since the other items are worth 0 to both, it doesn't matter who gets them. So, I should propose to take the remote and you should propose to take the pliers and boxes.} \\
\\ 
\texttt{But wait, the other player also wants to maximize the total score. If I propose to take the remote} \\
            }
        }
    }
    & \\ \\

    \theutterance \stepcounter{utterance}  
    & & & \multicolumn{4}{p{0.6\linewidth}}{
        \cellcolor[rgb]{0.9,0.9,0.9}{
            \makecell[{{p{\linewidth}}}]{
                \texttt{\tiny{[GM$\langle$P2]}}
                \texttt{{[}Proposal: 1 remote{]}} \\
            }
        }
    }
     \\ \\

    \theutterance \stepcounter{utterance}  
    & & & \multicolumn{2}{p{0.3\linewidth}}{
        \cellcolor[rgb]{0.9,0.9,0.9}{
            \makecell[{{p{\linewidth}}}]{
                \texttt{\tiny{[GM$|$GM]}}
                \texttt{proposal} \\
            }
        }
    }
    & & \\ \\

    \theutterance \stepcounter{utterance}  
    & & & \multicolumn{2}{p{0.3\linewidth}}{
        \cellcolor[rgb]{0.9,0.9,0.9}{
            \makecell[{{p{\linewidth}}}]{
                \texttt{\tiny{[GM$|$GM]}}
                [0, 0, 1]
            }
        }
    }
    & & \\ \\

    \theutterance \stepcounter{utterance}  
    & & \multicolumn{4}{p{0.6\linewidth}}{
        \cellcolor[rgb]{0.9,0.9,0.9}{
            \makecell[{{p{\linewidth}}}]{
                \texttt{\tiny{[P1$\langle$GM]}}
                \texttt{The other player has submitted a secret proposal. It is now your turn to submit a proposal of your own.} \\
            }
        }
    }
    & \\ \\

    \theutterance \stepcounter{utterance}  
    & \multicolumn{4}{p{0.6\linewidth}}{
        \cellcolor[rgb]{0.9,0.9,0.9}{
            \makecell[{{p{\linewidth}}}]{
                \texttt{\tiny{[P1$\rangle$GM]}}
                \texttt{{[}Proposal: 1 remote{]}} \\
            }
        }
    }
    & & \\ \\

    \theutterance \stepcounter{utterance}  
    & & & \multicolumn{2}{p{0.3\linewidth}}{
        \cellcolor[rgb]{0.9,0.9,0.9}{
            \makecell[{{p{\linewidth}}}]{
                \texttt{\tiny{[GM$|$GM]}}
                \texttt{proposal} \\
            }
        }
    }
    & & \\ \\

    \theutterance \stepcounter{utterance}  
    & & & \multicolumn{2}{p{0.3\linewidth}}{
        \cellcolor[rgb]{0.9,0.9,0.9}{
            \makecell[{{p{\linewidth}}}]{
                \texttt{\tiny{[GM$|$GM]}}
                [0, 0, 1]
            }
        }
    }
    & & \\ \\

    \theutterance \stepcounter{utterance}  
    & & & \multicolumn{2}{p{0.3\linewidth}}{
        \cellcolor[rgb]{0.9,0.9,0.9}{
            \makecell[{{p{\linewidth}}}]{
                \texttt{\tiny{[GM$|$GM]}}
                [[0, 0, 1], [0, 0, 1]]
            }
        }
    }
    & & \\ \\

    \theutterance \stepcounter{utterance}  
    & \multicolumn{4}{p{0.6\linewidth}}{
        \cellcolor[rgb]{0.95,0.95,0.95}{
            \makecell[{{p{\linewidth}}}]{
                \texttt{\tiny{[P2$\rightarrow$P1]}}
                \texttt{Post{-}episode Reflexion (not sent to GM):} \\
\texttt{New plan: I failed because I assumed the other player would accept a proposal where they receive items worth 0 to them, even though the total sum was maximized. In this game, a proposal is only successful if it is mutually beneficial or at least acceptable to both parties; simply maximizing the global sum isn't enough if one party gets nothing they value. Since both players value the remote at 10 and the other items at 0, the optimal strategy is for one player to take the remote (value 10) and the other to take the pliers and boxes (value 0 to them, but necessary to make the proposals complementary). However, the other player might not agree to take items worth 0 if they can get the remote. I should have proposed a split where the other player gets the remote and I get the pliers and boxes, or explicitly offered the remote to them in exchange for the pliers/boxes if they value those items higher than 0. Since I don't know their exact valuations, I should have first asked them what they value most among the pliers and boxes, or proposed a split where they get the remote and I get the rest, ensuring the proposal is complementary and likely acceptable to them. Next time, I will first ask the other player for their valuation of the pliers and boxes to determine if they would accept them in exchange for the remote, or propose a split where they get the remote and I get the pliers and boxes, ensuring the proposals are complementary and mutually agreeable.} \\
            }
        }
    }
    & & \\ \\

    \theutterance \stepcounter{utterance}  
    & & & \multicolumn{4}{p{0.6\linewidth}}{
        \cellcolor[rgb]{0.95,0.95,0.95}{
            \makecell[{{p{\linewidth}}}]{
                \texttt{\tiny{[P1$\rightarrow$P2]}}
                \texttt{Post{-}episode Reflexion (not sent to GM):} \\
\texttt{New plan: I failed because I submitted a proposal containing only the remote without offering any other items, which meant my proposal was not complementary to the other player's potential proposal (they would need to take the remote, leaving no items for me, or take nothing, failing the complementarity rule). Since the other player also values the remote at 10 and the other items at 0, they will likely propose to take the remote. To succeed, I must propose a bundle that includes the remote plus all the other items (pliers and boxes) so that if they take the remote, my proposal covers the rest, satisfying the complementarity condition. Next time, I will submit a proposal for "1 remote, 2 pliers, 3 boxes" immediately or after confirming their intent to take the remote, ensuring the union of our proposals covers all items.} \\
            }
        }
    }
     \\ \\

\end{supertabular}
}

\end{document}
